% Propietario conservado: /Users/ruben/Documents/ChatGPT/jueces y controles/REVISION_ARGUMENTAL_APERTURAS_CIERRES_20260915/II/manuscript_es/sections/definicion_barbero_rev06.tex
\subsection{Définition structurelle et fonction du paramètre de Barbero--Immirzi}
\label{ii:sec:ii-definicion-barbero}

\begin{definition}[Évaluation orientée de l'incidence hexadique--octadique]
Le paramètre \(\gamma_{\HMT}\) est l'évaluation conjuguée de la
chaîne fondamentale \(c_{12}^{+}\), complétée par sa composante
angulaire. Dans le domaine \(0<q_\pm<1\), sa définition
\eqref{ii:eq:gamma} s'exprime par
\begin{equation}
 \gamma_{\HMT}
 =\frac{\pi AC^*}{180}
  +[\mathscr R(c_{12}^{+})](q_-)
  -[\mathscr R(c_{12}^{+})](q_+).
 \label{ii:eq:ii-definicion-barbero-evaluacion}
\end{equation}
L'incidence marquée fournit les espaces d'événements ; la chaîne
orientée fixe leurs multiplicités ; les canaux de
\eqref{ii:eq:ii-canales} réalisent l'évaluation.
\end{definition}

La détermination utilise conjointement les cardinaux \(90,120\),
les douze contributions de secteur et le raccordement de retour. Le
coefficient du pôle \(1/24\) et la minimalité diophantienne vérifient
la compatibilité de ces poids avec la normalisation. Leur provenance
est conservée dans \(c_{12}^{+}\), avant l'évaluation des canaux. L'inversion
angulaire échange \(q_+\) et \(q_-\) et produit
\(\gamma_{\HMT}(A,-C^*)=-\gamma_{\HMT}(A,C^*)\) : la valeur
adimensionnelle transporte une orientation de l'incidence.

Dans la réalisation géométrique, ce même coefficient intervient dans
\eqref{ii:eq:holst} et dans l'incrément d'aire
\(\gamma_{\HMT}a_0\ell_P^2\). L'action de Holst utilise
\(\gamma_{\HMT}\ne0\) ; la grandeur aréale positive emploie
\(|\gamma_{\HMT}|\). La proposition~\ref{ii:prop:ii-costura-barbero}
explicite le changement coordonné d'orientation. Ainsi, la structure
du paramètre explique ce qu'il pondère dans l'incidence, ce qu'il conserve
lors du transport et comment il participe à la réponse torsionnelle et à
la résolution géométrique de l'information.
